\section*{Referencias complementarias de conservación y representación}
\addcontentsline{toc}{section}{Referencias complementarias de conservación y representación}
\begin{enumerate}
\item E. Noether, \emph{Invariante Variationsprobleme}, Nachrichten von der Gesellschaft der Wissenschaften zu Göttingen, Mathematisch-Physikalische Klasse (1918), 235--257. Traducción de M. A. Tavel: \href{https://arxiv.org/abs/physics/0503066}{Invariant Variation Problems}.
\item C. Weedbrook, S. Pirandola, R. García-Patrón, N. J. Cerf, T. C. Ralph, J. H. Shapiro y S. Lloyd, \emph{Gaussian Quantum Information}, Reviews of Modern Physics 84 (2012), 621--669. \href{https://arxiv.org/abs/1110.3234}{arXiv:1110.3234}.
\end{enumerate}
\begin{enumerate}[resume]
\item J. Williamson, \emph{On the Algebraic Problem Concerning the Normal Forms of Linear Dynamical Systems}, American Journal of Mathematics 58 (1936), 141--163. \href{https://doi.org/10.2307/2371062}{doi:10.2307/2371062}.
\item A. S. Holevo, M. Sohma y O. Hirota, \emph{Capacity of Quantum Gaussian Channels}, Physical Review A 59 (1999), 1820--1828. \href{https://doi.org/10.1103/PhysRevA.59.1820}{doi:10.1103/PhysRevA.59.1820}.
\end{enumerate}
